\begin{abstract} 
Real user studies are important for understanding how people interact with systems under test or already deployed. In practice, however, they are often costly, time-consuming, and difficult to scale. To address these challenges, we introduce \method, a persona-based user simulation framework that approximates real-user behavior across diverse interactive settings. \method connects simulated users drawn from existing persona datasets to task-specific application interfaces and collects the interaction trajectories and outcomes. 
% These simulated users are instantiated with diverse, real-life persona profiles that include backgrounds, personalities, goals, and other relevant attributes. 
\method  provides a plug-and-play evaluation workflow in which the application being evaluated can be easily changed. In this demo, we present \method on three forms of interactive applications: surveys, 
% where simulated users answer structured questionnaires; 
chatbots, 
% where they complete multi-turn conversations with a chatbot;
and web applications.
% where they navigate web applications. 
Together, these examples show that \method can support repeatable, parallelizable, and scalable evaluation across different interaction settings, while producing user-oriented feedback and task-specific behavior.
\end{abstract}

% （rosie modified version to emphasis our selling point)
% \begin{abstract} 
% Real-user studies are essential for evaluating interactive AI systems, but they are often expensive, slow, and difficult to scale during iterative development. We present PersonaEval, a unified persona-grounded user simulation framework for evaluating interactive applications across diverse domains. PersonaEval instantiates simulated users from existing persona datasets containing attributes such as background, personality, goals, and preferences, and connects them to task-specific application interfaces to execute realistic interaction workflows. Its plug-and-play architecture enables new applications to be integrated with minimal engineering effort while collecting interaction trajectories and task outcomes for evaluation. We demonstrate PersonaEval on three representative interaction settings: structured surveys, conversational recommendation agents, and interactive web applications. These examples illustrate how PersonaEval enables scalable, repeatable, and low-cost user-centered evaluation while providing rich interaction traces and persona-specific analyses across heterogeneous applications.
% \end{abstract}